% GENERATED FILE. Edit canonical lessons, then run npm run generate:books.
% canonical-source-hash: fnv1a64:8d70dfaf5f9ad3e9
% canonical-lessons: MR-C229-goli, MR-C229-malam, MR-C229-patti, MR-C229-rugnavahika, MR-C229-ahar

\chapter{Tablet, Ointment, Bandage, Ambulance, Diet}
\label{ch:mr-tablet-ointment-bandage-ambulance-diet}
\hlchaptermodality{229}

\begin{chapteropening}
\textbf{By the end of this chapter:} \emph{I can say a tablet, an ointment, a bandage, an ambulance and a diet.}

Complete the last lesson of chapter 229: I can say a tablet, an ointment, a bandage, an ambulance and a diet.
\end{chapteropening}

\section[\texorpdfstring{\mr{गोळी} (goḷī)}{गोळी (goḷī)}]{\mr{गोळी} (goḷī) — a tablet, a pill}

\label{lesson:MR-C229-goli}



\begin{quote}
Before the new one: say the Marathi for to defend, then the Marathi for to disappear.
\end{quote}

\subsection*{You'll want to know: \mr{गोळी}}
\textbf{\mr{गोळी}} — \emph{goḷī} — \textquotedblleft{}a tablet, a pill\textquotedblright{}.

\begin{grammarlens}[title={What you've built}]
One more word for this chapter.
\end{grammarlens}

\subsection*{Your turn}
\begin{itemize}
\raggedright
  \item \emph{Say it:} \emph{goḷī}
  \item \emph{Say it:} \emph{goḷī}, once more
  \item \emph{Say it:} say \emph{goḷī}
  \item \emph{Recall:} say the Marathi for to defend, then the Marathi for to disappear, then say \emph{goḷī} again
\end{itemize}

\begin{tcolorbox}[breakable,colback=teal!4,colframe=teal!35!black,title={Before you move on}]
What does \mr{गोळी} mean? (\textquotedblleft{}A tablet, a pill\textquotedblright{}.) Say it once more.
\end{tcolorbox}
\section[\texorpdfstring{\mr{मलम} (malam)}{मलम (malam)}]{\mr{मलम} (malam) — an ointment}

\label{lesson:MR-C229-malam}



\begin{quote}
Before the new one: say the Marathi for to disappear, then the Marathi for a tablet.
\end{quote}

\subsection*{You'll want to know: \mr{मलम}}
\textbf{\mr{मलम}} — \emph{malam} — \textquotedblleft{}an ointment\textquotedblright{}.

\begin{grammarlens}[title={What you've built}]
One more word for this chapter.
\end{grammarlens}

\subsection*{Your turn}
\begin{itemize}
\raggedright
  \item \emph{Say it:} \emph{malam}
  \item \emph{Say it:} \emph{malam}, once more
  \item \emph{Say it:} say \emph{malam}
  \item \emph{Recall:} say the Marathi for to disappear, then the Marathi for a tablet, then say \emph{malam} again
\end{itemize}

\begin{tcolorbox}[breakable,colback=teal!4,colframe=teal!35!black,title={Before you move on}]
What does \mr{मलम} mean? (\textquotedblleft{}An ointment\textquotedblright{}.) Say it once more.
\end{tcolorbox}
\section[\texorpdfstring{\mr{पट्टी} (paṭṭī)}{पट्टी (paṭṭī)}]{\mr{पट्टी} (paṭṭī) — a bandage; a strip}

\label{lesson:MR-C229-patti}



\begin{quote}
Before the new one: say the Marathi for a tablet, then the Marathi for an ointment.
\end{quote}

\subsection*{You'll want to know: \mr{पट्टी}}
\textbf{\mr{पट्टी}} — \emph{paṭṭī} — \textquotedblleft{}a bandage; a strip\textquotedblright{}.

\begin{grammarlens}[title={What you've built}]
One more word for this chapter.
\end{grammarlens}

\subsection*{Your turn}
\begin{itemize}
\raggedright
  \item \emph{Say it:} \emph{paṭṭī}
  \item \emph{Say it:} \emph{paṭṭī}, once more
  \item \emph{Say it:} say \emph{paṭṭī}
  \item \emph{Recall:} say the Marathi for a tablet, then the Marathi for an ointment, then say \emph{paṭṭī} again
\end{itemize}

\begin{tcolorbox}[breakable,colback=teal!4,colframe=teal!35!black,title={Before you move on}]
What does \mr{पट्टी} mean? (\textquotedblleft{}A bandage; a strip\textquotedblright{}.) Say it once more.
\end{tcolorbox}
\section[\texorpdfstring{\mr{रुग्णवाहिका} (rugṇavāhikā)}{रुग्णवाहिका (rugṇavāhikā)}]{\mr{रुग्णवाहिका} (rugṇavāhikā) — an ambulance}

\label{lesson:MR-C229-rugnavahika}



\begin{quote}
Before the new one: say the Marathi for an ointment, then the Marathi for a bandage; a strip.
\end{quote}

\subsection*{You'll want to know: \mr{रुग्णवाहिका}}
\textbf{\mr{रुग्णवाहिका}} — \emph{rugṇavāhikā} — \textquotedblleft{}an ambulance\textquotedblright{}.

\begin{grammarlens}[title={What you've built}]
One more word for this chapter.
\end{grammarlens}

\subsection*{Your turn}
\begin{itemize}
\raggedright
  \item \emph{Say it:} \emph{rugṇavāhikā}
  \item \emph{Say it:} \emph{rugṇavāhikā}, once more
  \item \emph{Say it:} say \emph{rugṇavāhikā}
  \item \emph{Recall:} say the Marathi for an ointment, then the Marathi for a bandage; a strip, then say \emph{rugṇavāhikā} again
\end{itemize}

\begin{tcolorbox}[breakable,colback=teal!4,colframe=teal!35!black,title={Before you move on}]
What does \mr{रुग्णवाहिका} mean? (\textquotedblleft{}An ambulance\textquotedblright{}.) Say it once more.
\end{tcolorbox}
\section[\texorpdfstring{\mr{आहार} (āhār)}{आहार (āhār)}]{\mr{आहार} (āhār) — a diet}

\label{lesson:MR-C229-ahar}



\begin{quote}
Before the new one: say the Marathi for a bandage; a strip, then the Marathi for an ambulance.
\end{quote}

\subsection*{You'll want to know: \mr{आहार}}
\textbf{\mr{आहार}} — \emph{āhār} — \textquotedblleft{}a diet\textquotedblright{}.

\begin{grammarlens}[title={What you've built}]
One more word for this chapter.
\end{grammarlens}

\subsection*{Your turn}
\begin{itemize}
\raggedright
  \item \emph{Say it:} \emph{āhār}
  \item \emph{Say it:} \emph{āhār}, once more
  \item \emph{Say it:} say \emph{āhār}
  \item \emph{Recall:} say the Marathi for a bandage; a strip, then the Marathi for an ambulance, then say \emph{āhār} again
\end{itemize}

\begin{tcolorbox}[breakable,colback=teal!4,colframe=teal!35!black,title={Before you move on}]
What does \mr{आहार} mean? (\textquotedblleft{}A diet\textquotedblright{}.) Say it once more.
\end{tcolorbox}
